\chapter{Production Agent Systems}
\label{ch:production}

\epigraph{``A demo agent runs once. A production agent runs until it fails---and then it must recover.''}{}

\learninggoal{
	\item Design a production agent architecture with observability, reliability, and scalability.
	\item Implement caching, rate limiting, and cost management.
	\item Understand the operational challenges unique to agent systems.
	\item Compare deployment patterns for synchronous and asynchronous agents.
}

\section{Production Architecture}

A production agent system extends the core loop with infrastructure for reliability, observability, and scale.

\begin{lstlisting}[caption=Production agent architecture]
                    +-----------+
                    |  Client   |
                    +-----+-----+
                          |
                    +-----v-----+
                    |  API GW   |  <- Auth, rate limiting, routing
                    +-----+-----+
                          |
                    +-----v-----+
                    | Scheduler |  <- Queue, prioritise, retry
                    +-----+-----+
                          |
                    +-----v------+
                    | Agent      |  <- Core loop (Ch. 2)
                    | Executor   |
                    +-----+------+
                    |     |      |
              +-----v--+--v---  --v-----+
              | Tool   | Memory | Cache |
              | Exec   | Store  | Layer |
              +--------+-------+-------+
                    |     |      |
              +-----v-----v------v-----+
              | Observability Stack    |
              | (logging, metrics,     |
              |  tracing, alerting)    |
              +------------------------+
\end{lstlisting}

\section{Request Lifecycle}

\begin{lstlisting}[caption=Production request flow]
  1. Receive: API gateway authenticates, rate-limits, routes
  2. Queue: Scheduler adds to priority queue
  3. Dispatch: Agent executor picks up (may wait for capacity)
  4. Process: Agent loop executes (checkpointed)
  5. Respond: Result returned, resources freed
  6. Audit: Full trajectory logged to data warehouse
\end{lstlisting}

\subsection{Synchronous vs. Asynchronous}

\begin{longtable}{p{3cm}p{5cm}p{5cm}}
	\toprule
	\textbf{Property} & \textbf{Synchronous}         & \textbf{Asynchronous}              \\
	\midrule
	API style         & HTTP POST, wait for response & Submit job, poll for result        \\
	Latency           & Low (seconds)                & High (minutes--hours)              \\
	Throughput        & Lower (blocking)             & Higher (queue-based)               \\
	Complexity        & Simple                       & Webhook or polling loop            \\
	Best for          & Simple QA, quick actions     & Complex tasks, long-running agents \\
	Error handling    & Return HTTP error            & Retry queue, dead-letter queue     \\
	\bottomrule
\end{longtable}

\section{Caching}

Agents issue many repeated \llm\ calls and tool calls. Caching reduces latency and cost.

\subsection{Semantic Caching}

Two semantically identical queries should not require two \llm\ calls:

\begin{lstlisting}[caption=Semantic cache for agent responses]
  class SemanticCache:
      def __init__(self, threshold=0.95):
          self.cache = []  // [{embedding, response}]
          self.threshold = threshold

      def get(self, query):
          query_emb = embed(query)
          for entry in self.cache:
              sim = cosine_similarity(query_emb, entry.embedding)
              if sim >= self.threshold:
                  return entry.response
          return None

      def set(self, query, response):
          self.cache.append({
              embedding: embed(query),
              response: response
          })
\end{lstlisting}

\subsection{Tool Result Caching}

Idempotent tool calls with the same arguments can be cached:

\begin{lstlisting}[caption=Tool result cache]
  def cached_tool_call(name, args, ttl=300):
      key = f"{name}:{json.dumps(args, sort_keys=True)}"
      cached = redis.get(key)
      if cached:
          return json.loads(cached)

      result = execute_tool(name, args)
      if tool_is_idempotent(name):
          redis.setex(key, ttl, json.dumps(result))
      return result
\end{lstlisting}

\subsection{KV Cache Reuse}

When the system prompt is identical across requests, the KV cache for the prefix can be shared:

\begin{itemize}
	\item All requests that share a system prompt reuse the same KV cache prefix.
	\item Common prefixes (tool definitions, few-shot examples) are computed once.
	\item Prefix caching (see the companion volume on serving) reduces TTFT by 30--50\%.
\end{itemize}

\section{Reliability Patterns}

\subsection{Retry with Backoff}

Transient failures (rate limits, network errors) should be retried:

\begin{lstlisting}[caption=Exponential backoff for tool calls]
  async def retry_with_backoff(fn, max_retries=3):
      for attempt in range(max_retries):
          try:
              return await fn()
          except TransientError as e:
              if attempt == max_retries - 1:
                  raise
              wait = min(2 ** attempt * 0.1, 5.0)  # 0.1, 0.2, 0.4, ...
              await asyncio.sleep(wait)
\end{lstlisting}

\subsection{Circuit Breaker}

Prevent cascading failures when a downstream service is degraded:

\begin{lstlisting}[caption=Circuit breaker pattern]
  class CircuitBreaker:
      states = ("closed", "open", "half-open")

      def __init__(self, threshold=5, recovery_time=30):
          self.state = "closed"
          self.failure_count = 0
          self.threshold = threshold
          self.recovery_time = recovery_time
          self.last_failure = None

      async def call(self, fn):
          if self.state == "open":
              if time() - self.last_failure > self.recovery_time:
                  self.state = "half-open"
              else:
                  raise CircuitOpenError()

          try:
              result = await fn()
              if self.state == "half-open":
                  self.state = "closed"
                  self.failure_count = 0
              return result
          except Exception:
              self.failure_count += 1
              self.last_failure = time()
              if self.failure_count >= self.threshold:
                  self.state = "open"
              raise
\end{lstlisting}

\subsection{Dead-Letter Queue}

Requests that fail after all retries go to a dead-letter queue for manual inspection:

\begin{lstlisting}[caption=Dead-letter queue processing]
  # When all retries are exhausted:
  dead_letter_queue.send({
      "request": original_request,
      "trajectory": agent_trajectory,
      "error": final_error,
      "attempts": 3,
      "timestamp": now()
  })

  # Manual review process:
  for entry in dead_letter_queue:
      review = human_review(entry)
      if review.action == "retry":
          agent_executor.submit(review.modified_request)
      elif review.action == "ignore":
          entry.archive()
\end{lstlisting}

\section{Cost Management}

Agent costs can spiral. Tracking and budgeting is essential:

\begin{longtable}{p{3cm}p{4cm}p{4cm}p{3cm}}
	\toprule
	\textbf{Cost component} & \textbf{Driver}       & \textbf{Mitigation}      & \textbf{Impact} \\
	\midrule
	LLM inference           & Tokens per trajectory & Caching, shorter prompts & 40--60\%        \\
	Tool execution          & API calls, compute    & Cache idempotent calls   & 10--20\%        \\
	Embedding               & Memory retrieval      & Cache embeddings         & 5--10\%         \\
	Human review            & Approval gates        & Reduce unnecessary gates & 10--30\%        \\  % chktex 27
	Infrastructure          & Serving, storage      & Autoscaling              & 10--20\%        \\
	\bottomrule
\end{longtable}

\begin{principle}[Cost Awareness]
	The agent system should track its own cost per session and make cost-visible decisions: ``This search costs \$0.02; I will perform it.'' For high-cost operations, the agent should ask: ``Is it worth it?''
\end{principle}

\section{Observability}

\subsection{Logging}

Every agent turn is logged:

\begin{lstlisting}[caption=Agent log entry]
  {
    "session_id": "abc-123",
    "turn": 5,
    "timestamp": "2025-05-24T12:34:56Z",
    "duration_ms": 2340,
    "llm_call": {
      "model": "gpt-4o",
      "prompt_tokens": 4520,
      "completion_tokens": 312,
      "cost": 0.023
    },
    "tool_call": {
      "name": "search",
      "args_summary": "q=population+Tokyo",
      "duration_ms": 890,
      "success": true
    },
    "decision": {
      "type": "tool_call",
      "confidence": 0.85
    }
  }
\end{lstlisting}

\subsection{Metrics}

Key performance indicators for production agents:

\begin{itemize}
	\item \textbf{P50/P99 latency:} end-to-end request duration.
	\item \textbf{Success rate:} fraction of requests completed without error.
	\item \textbf{Intervention rate:} fraction of requests requiring human approval.
	\item \textbf{Cost per request:} average \$ per completed request.
	\item \textbf{LLM call count:} average number of \llm\ calls per request.
	\item \textbf{Tool failure rate:} fraction of tool calls that error.
	\item \textbf{Cache hit rate:} fraction of queries served from cache.
\end{itemize}

\subsection{A/B Testing}

Changes to agent prompts or tool sets should be evaluated via A/B testing:

\begin{lstlisting}[caption=A/B test configuration]
  ABTest ::= {
    experiment_id: UUID,
    variants: {
      "control": { system_prompt: "..." },
      "treatment": { system_prompt: "...", tools: [...] }
    },
    traffic_split: { control: 0.5, treatment: 0.5 },
    metrics: ["success_rate", "cost_per_task", "latency"],
    minimum_samples: 1000,
    duration: "7d"
  }
\end{lstlisting}

\section{Scaling}

\subsection{Horizontal Scaling}

Agent executors are stateless (state lives in the message store), enabling horizontal scaling:

\begin{itemize}
	\item Add more executor instances behind the scheduler.
	\item Each executor processes one request at a time (LLM calls are sequential within a session).
	\item Executors share nothing except the cache and memory store.
\end{itemize}

\subsection{Limits}

The bottleneck is typically \llm\ API rate limits:

\begin{equation}
	\text{Max throughput} = \min\left(\frac{\text{\# executors}}{\text{avg time per request}}, \text{LLM rate limit}\right)
\end{equation}

For a 30-second average request and 100 executors: max throughput $\approx 200$ requests/minute, unless the \llm\ API rate limit is lower.

\section{Exercises}

\begin{exercise}
	Design a cost-tracking system for an agent that uses GPT-4o (\$10/M input tokens, \$30/M output tokens) and makes 10 web search calls (\$0.01 each) per request. What is the cost ceiling per request? How would you enforce it?
\end{exercise}

\begin{exercise}
	Implement a semantic cache for an agent that answers customer support questions. Measure the cache hit rate on 10,000 historical queries. What similarity threshold gives the best trade-off between hit rate and correctness?
\end{exercise}

\begin{exercise}
	Design a circuit breaker for an agent's database tool. The database is slow. At what point should the circuit breaker open? What happens to in-flight requests? What is the recovery strategy?
\end{exercise}

\begin{exercise}
	Compare synchronous and asynchronous deployment for a code-review agent. A code review takes 2--5 minutes. Which deployment model is better? Design the API contract for your chosen model.
\end{exercise}
